\begin{proposition}[Simultaneous Phase--Depth Extraction via Scan Spectral Tail Ratios: Locating the Dominant Defect without Full Inversion]\label{prop:xi-scan-tail-ratio-dominant-defect}
Under the exponential spectral representation of Proposition \ref{con:xi-2kappa-minimal-complete-samples}, let
\[
u_n=\sum_{j=1}^{\kappa}c_j\,r_j^{n}\qquad(n\ge 0),
\qquad
r_j=e^{-(1+\delta_j)\Delta}e^{-\mathrm{i}\gamma_j\Delta},
\qquad c_j>0.
\]
Assume strict modulus separation: there exists an ordering such that
\[
|r_1|>|r_2|\ge\cdots\ge|r_\kappa|.
\]
Then the limit
\[
\boxed{\;\lim_{n\to\infty}\frac{u_{n+1}}{u_n}=r_1\;}
\]
exists, and convergence is exponentially fast: there exists a constant \(C>0\) such that
\[
\left|\frac{u_{n+1}}{u_n}-r_1\right|
\ \le\ C\left(\frac{|r_2|}{|r_1|}\right)^{n}\qquad(n\to\infty).
\]
Therefore, the depth and height congruence class of the dominant defect can be simultaneously recovered from the tail ratio:
\[
\boxed{\;
\delta_{\min}=\frac{-\log|r_1|}{\Delta}-1,
\qquad
\gamma_1\equiv-\frac{1}{\Delta}\arg(r_1)\ \Bigl(\bmod\ \frac{2\pi}{\Delta}\Bigr)\; }.
\]
\end{proposition}
\begin{proof}
Write \(u_n=c_1r_1^n(1+\varepsilon_n)\), where
\(
\varepsilon_n:=\sum_{j\ge 2}\frac{c_j}{c_1}\bigl(\frac{r_j}{r_1}\bigr)^n
\).
From \(|r_j/r_1|\le |r_2/r_1|<1\) we obtain \(\varepsilon_n=O((|r_2|/|r_1|)^n)\), so
\[
\frac{u_{n+1}}{u_n}
=r_1\,\frac{1+\varepsilon_{n+1}}{1+\varepsilon_n}
=r_1\Bigl(1+O\bigl((|r_2|/|r_1|)^n\bigr)\Bigr).
\]
The final two formulas follow from \(|r_1|=e^{-(1+\delta_1)\Delta}\) and \(\arg(r_1)\equiv-\gamma_1\Delta\ (\bmod\ 2\pi)\). \qed
\end{proof}

\begin{proposition}[Successive Peeling Theorem under Strict Modulus Separation: Tail-Limit-Driven Spectral Decomposition Closes Exactly in Finitely Many Steps]\label{prop:xi-scan-successive-peeling-closure}
Under the assumptions of Proposition \ref{prop:xi-scan-tail-ratio-dominant-defect}, let
\[
r_1:=\lim_{n\to\infty}\frac{u_{n+1}}{u_n},
\qquad
c_1:=\lim_{n\to\infty}\frac{u_{n}}{r_1^{n}}.
\]
Define the peeled sequence
\[
u_n^{(1)}:=u_n-c_1r_1^n.
\]
Then \(u_n^{(1)}=\sum_{j=2}^{\kappa}c_jr_j^n\), and the same type of modulus separation still holds.
Iterating \(\kappa\) times yields exact closure in finitely many steps: there exists a family \(\{(c_j,r_j)\}_{j=1}^{\kappa}\) uniquely determined by successive tail limits such that
\[
u_n=\sum_{j=1}^{\kappa}c_jr_j^n\qquad(\forall n\ge 0).
\]
Therefore, in the regime subclass with strict modulus separation, the complete defect spectrum can be exactly recovered by finitely many steps of tail-limit-driven spectral peeling without requiring a ``global one-shot inversion''; the unavoidable ill-posedness occurs only at modulus collision \(|r_j|\to|r_k|\), consistent with the ``collision implies recovery'' ill-posedness mechanism of Proposition \ref{con:xi-collision-illposedness-sep-blowup}.
\end{proposition}
\begin{proof}
By modulus separation, Proposition \ref{prop:xi-scan-tail-ratio-dominant-defect} yields \(r_1\).
From \(u_n=c_1r_1^n+O(|r_2|^n)\), the limit \(c_1=\lim_{n\to\infty}u_n/r_1^n\) exists.
Substituting into the definition yields the peeling identity \(u_n^{(1)}=\sum_{j\ge 2}c_jr_j^n\).
Repeating the above argument for \(u^{(1)}\) yields \(r_2,c_2\), and iterating to the \(\kappa\)-th step recovers the full spectrum with uniqueness. \qed
\end{proof}

\endinput
